% !TEX root = ../paper.tex
\section{Preliminaries} \label{sec:preliminaries}

In this section, we introduce concepts and 
notation that we use for the rest of the paper.
Let $F = (X, \cC)$ be a $k$-CNF with variable set $X$ and clause set $\cC$. We view the satisfying assignments of $F$ as 
subsets of variables which are set to 1.

\begin{definition}[Transversals]
A set $S\subseteq X$ is a transversal for $F$ if the assignment that sets
the variables in $S$ to $1$ and the variables in 
$X\setminus S$ to $0$ is a satisfying assignment of $F$.
The size of a transversal $S$ is defined as $|S|$.
\end{definition}

We say $S$ is a minimal transversal for $F$ if no subset of $S$ is a transversal. We call the corresponding satisfying asssignment a minimal satisfying assignment.

Our use of transversals for discussing satisfying assignments is
motivated by the notion of transversals in hypergraphs.
We view a monotone 
$k$-CNF 
(where all literals in every clause are positive) 
as a $k$-graph where every hyperedge has size at most $k$. 
This leads to a 1-1 correspondence
between the transversals of a monotone
$k$-CNF 
and those  of the corresponding hypergraph.
%We note that for a non-monotone $k$-CNF $F = (X, \cC)$, 
%there could exist 
%$A\subseteq B\subseteq C\subseteq X$ such that $A$ is a transversal, $B$ is not a transversal, and $C$ is a transversal.

In this paper, we are primarily interested in 
minimum-size transversals.

\begin{definition}[Transversal number]
%We say $S\subseteq X$ is a minimum-size transversal 
%if there are no transversals of size
%for all $T\subsetneq S$, $T$ is not a transversal.
%We say that $S$ is a $t$-transversal 
%if $S$ is a minimal transversal
%and $|S|=t$.
For a satisfiable $k$-CNF $F$, we define \emph{transversal number} $\trn(F)$ to be the cardinality of the minimum-size transversal of $F$.  
We use $\trans(F)$ to denote the set of all minimum-size
transversals of $F$
and $\ntrn(F)$ to denote the cardinality of 
$\trans(F)$.
% number of $t$-transversals of $F$.
%For convenience, we use 
%$\trans(F)$ to denote $\trans(F)_{\trn(F)}$ and
%$\ntrn(F)$ to denote $\ntrn(F)_{\trn(F)}$.
\end{definition}


Let $t\in [n]$ and $\alpha\in \zo^n$ be such 
that every satisfying assignment of $F$ is at a distance
of at least $t$ from $\alpha$. 
We reduce the $\enumballsat{k}{t}$  problem for $F$
from the initial assignment $\alpha$
to the $\enumz{k}$  problem:
enumerate all 
minimum-size transversals of a $k$-CNF\footnote{This problem has been previously considered by, e.g. \cite{FominGLS19} with a different name {\sc Min-Ones $k$-SAT} and a non-trivial algorithm is given independent of $\trn(F)$. Here we focus on running times with fine dependence on $\trn(F)$.}.
%$t$ starting from the all-0 assignment where we assume that $F$ contains
%no transversals of Hamming distance 
%less than $t$ from the all-0 assignment.
%We remark that  the general 
%$\enumballsat{k}{t}$ can be reduced to $\enumz{k}{t}$.
%Let $\alpha$ be the initial assignment 
%for the $\enumballsat{k}{t}$ problem with input $F$.
Indeed, let $G$ be the $k$-CNF formula (over 
the variable set $X$) where its
clause set is obtained from that of
$F$ by the following replacement of literals:  
For each variable $v\in X$, if $\alpha$ sets $v$ to $1$, 
then we swap occurrences of the positive and the negative 
literals corresponding to $v$ in $\cC$. 
Otherwise, if $\alpha$ sets $v$ to $0$,  we leave the corresponding literals
as they are.
$G$ is also a $k$-CNF and 
for all $y\in \zo^n$, $G(y) = F(y\oplus \alpha)$. Clearly, 
there exists a transversal within distance $t$ 
from $0^n$ in $G$ if and only if there exists a transversal 
within distance $t$ from $\alpha$ in $F$. 

We now prove the following useful proposition that
allows us to assume all clauses of $F$ have width exactly $k$.

\begin{proposition}
\label{prop:monotone-k-cnf-uniform}
For every $k$-CNF $F = (X, \cC)$ with $\trn(F) \le n-k$,
there exists a $k$-CNF $G$ where every clause has width exactly $k$
and $\trn(G) = \trn(F)$
and the set of transversals of $G$ includes the
set of transversals of $F$.
Furthermore, if $F$ is monotone, $G$ is monotone.
\end{proposition}

\begin{proof}
% Let $F=(X,\cC)$ be a $k$-CNF where $\trn(F)\leq n-k$.
Assume that $F$ has a clause $C$ of width $1\leq k'<k$.
Let $F'$ be the formula obtained from $F$ by removing the clause $C$
and adding a clause  $C' = C\cup C'' $ for
each $S\subseteq X\setminus C, |S|= k-k'$, where $C''$ is a monotone clause over variables in $S$.
The proposition follows since every transversal of
$F$ is a transversal of $F'$ and any transversal
of $F'$ which is not a transversal of $F$ must have size at least $n-k+1$.
\end{proof}
